\documentclass[a4paper,11pt]{article}

\usepackage[T1]{fontenc}
\usepackage[utf8]{inputenc}
\usepackage[italian]{babel}
\usepackage{lmodern}
\usepackage{microtype}
\usepackage{geometry}
\usepackage{booktabs}
\usepackage{siunitx}
\usepackage{caption}
\usepackage{xcolor}
\usepackage{amsmath}
\usepackage{amssymb}
\usepackage{array}
\usepackage{url}

\geometry{margin=2.5cm}

\sisetup{
  output-decimal-marker = {,},
  group-separator       = {\,},
  group-minimum-digits  = 4,
  group-digits          = integer,
  detect-all
}

\title{Risultati della simulazione\\ \textit{MEV Insurance}}
\author{Alessandro}
\date{\today}

\begin{document}
\maketitle

\section{Configurazione}
La simulazione è stata eseguita con 7 oracoli, 1 bot attaccante (sandwich) e
10 trader, di cui 5 dotati di copertura \texttt{HIGH} e 5 privi di copertura
(\textit{naked}). Il pool AMM è stato inizializzato con \num{5000}~MEVI e
\num{5000}~USDC; il pool assicurativo è partito da \num{50}~MEVI. Il numero
complessivo di swap eseguiti è pari a \num{13212}. Lo scenario corrisponde al
caso~2.1 (\textit{5 insured vs 5 naked --- attack rate 10\,\%}).

La sessione ha avuto inizio alle ore 01:18 e si è conclusa alle ore 11:47, per
una durata complessiva di circa dieci ore e ventinove minuti. Poiché ogni minuto
reale corrispondeva a un giorno simulato, il sistema ha attraversato
\num{629}~giorni di attività compressi nel tempo reale. Ogni trader eseguiva in
media circa tre swap al giorno con valori di transazione compresi tra \num{50}
e \num{150}~USDC, generando un flusso continuo e realistico di operazioni sul
pool AMM per l'intera durata della simulazione.

\section{Attività dell'attaccante}
Il bot ha attaccato \num{1129} swap su \num{13212} totali, pari all'
\num{8,55}\,\%, sottraendo ai trader complessivamente \num{3979,11}~MEVI.
La perdita media per swap attaccato è dunque
\[
  \frac{3979{,}11}{1129} \;=\; \num{3,52}\ \text{MEVI}.
\]
Considerando il volume totale degli swap pari a \num{983515,22}~MEVI,
corrispondente ad uno swap medio di \num{74,44}~MEVI, la perdita media per
attacco rappresenta il \num{4,73}\,\% del valore della singola transazione.

\section{Riepilogo per trader}
La Tabella~\ref{tab:trader} riporta, per ciascun trader, il volume
swappato in MEVI, l'importo complessivamente sottratto dal bot, il premio
versato e il rimborso ricevuto. Tutti gli importi sono espressi in MEVI.

\begin{table}[h]
\centering
\caption{Sommario per singolo trader.}
\label{tab:trader}
\small
\begin{tabular}{lc S[table-format=9.4] S[table-format=4.4] S[table-format=6.4] S[table-format=6.4]}
\toprule
\textbf{Trader} & \textbf{Cop.}
  & {\textbf{Volume}} & {\textbf{Attaccato}}
  & {\textbf{Premio}} & {\textbf{Rimborso}} \\
\midrule
\texttt{0x14dc7996\dots} & ---  & 139251.4659 & 617.2902 & {---}      & {---}      \\
\texttt{0x23618e81\dots} & ---  &  83922.6746 & 255.8103 & {---}      & {---}      \\
\texttt{0x71be63f3\dots} & ---  & 139786.3399 & 627.5147 & {---}      & {---}      \\
\texttt{0xa0ee7a14\dots} & ---  & 138079.1204 & 490.7035 & {---}      & {---}      \\
\texttt{0xbcd4042d\dots} & ---  &  84268.6757 & 352.5365 & {---}      & {---}      \\
\midrule
\texttt{0x15d34aaf\dots} & HIGH &  71525.0061 & 251.2716 &  274.0831  &  239.9213  \\
\texttt{0x3c44cddd\dots} & HIGH &  70772.8040 & 256.1291 &  271.1824  &  243.4103  \\
\texttt{0x90f79bf6\dots} & HIGH &  71526.9068 & 321.8455 &  277.1605  &  300.2928  \\
\texttt{0x976ea740\dots} & HIGH &  72127.3627 & 303.3079 &  275.9357  &  280.0263  \\
\texttt{0x9965507d\dots} & HIGH & 112254.8627 & 502.7021 &  432.4858  &  436.5798  \\
\midrule
\multicolumn{2}{l}{\textbf{Totale non assicurati}}
       & 585308.2765 & 2343.8552 & {---}      & {---}      \\
\multicolumn{2}{l}{\textbf{Totale assicurati}}
       & 398206.9423 & 1635.2562 & 1530.8475  & 1500.2305  \\
\multicolumn{2}{l}{\textbf{Totale}}
       & 983515.2188 & 3979.1114 & 1530.8475  & 1500.2305  \\
\bottomrule
\end{tabular}
\end{table}

\section{Attaccato e rimborso: perché differiscono}
Nei trader assicurati si osserva una sistematica differenza tra l'importo
\emph{attaccato} --- cioè la perdita effettivamente subita durante lo swap ---
e il \emph{rimborso} erogato dal protocollo. Nel complesso, su
\num{1635,26}~MEVI dichiarati come perdita, ne sono stati rimborsati
\num{1500,23}; la differenza di \num{135,03}~MEVI (pari all'\num{8,26}\,\% dei
\emph{claim}) è stata rigettata dagli oracoli.

La causa risiede nel meccanismo di verifica: il protocollo accetta un rimborso
soltanto se il pattern dell'attacco sandwich registrato on-chain soddisfa i
criteri previsti dal contratto. Quando un trader invia un \emph{claim} con un
pattern non conforme --- ad esempio a causa di una stima errata degli indici
di prezzo pre/post-attacco o di una struttura della transazione non
riconosciuta come sandwich valido --- gli oracoli rigettano la richiesta. Il
trader subisce comunque la perdita, ma non riceve rimborso per quella
transazione.

Questo meccanismo è \textbf{intenzionale}: tutela la riserva del pool da
rimborsi fraudolenti o basati su errori lato client, garantendo che solo le
perdite verificabili vengano indennizzate.

\section{Effetto economico della copertura}
I trader assicurati hanno versato premi per \num{1530,85}~MEVI a fronte di
un volume swappato di \num{398206,94}~MEVI: l'incidenza media del premio sul
volume coperto è quindi pari allo \num{0,38}\,\%. I rimborsi complessivamente
erogati ammontano a \num{1500,23}~MEVI, da cui un \emph{loss ratio} del
\[
  \frac{1500{,}2305}{1530{,}8475} \;=\; \num{98,00}\,\%.
\]
Il pool assicurativo, partito da \num{50}~MEVI, ha registrato una crescita fino
a \num{80,62}~MEVI al termine della simulazione ($\Delta = +\num{30,62}$~MEVI).
Un \emph{loss ratio} del \num{98}\,\% indica che i premi raccolti coprono
abbondantemente i rimborsi erogati, lasciando un margine del \num{2}\,\% a
favore del pool. Dal lato del protocollo, questo margine è destinato a coprire
i costi operativi --- in particolare le ricompense degli oracoli --- e a far
crescere la riserva, garantendo la sostenibilità del meccanismo nelle condizioni
simulate.

\section{Confronto sul singolo swap}
Per illustrare l'utilità economica della copertura sul singolo trade, si
confrontano due transazioni colpite dal bot ed eseguite a poco più di un minuto
di distanza l'una dall'altra (Tabella~\ref{tab:confronto}).

\begin{table}[h]
\centering
\caption{Confronto tra una swap assicurata ed una non assicurata.}
\label{tab:confronto}
\small
\begin{tabular}{l S[table-format=3.2] S[table-format=3.4] S[table-format=1.4] S[table-format=1.4] S[table-format=1.4]}
\toprule
                 & {\textbf{Input (USDC)}} & {\textbf{Output (MEVI)}}
                 & {\textbf{Premio}} & {\textbf{Perdita}} & {\textbf{Rimborso}} \\
\midrule
Assicurata,    \texttt{01:26:13} &  77.25 &  56.3657 & 0.2009 & 1.9764 & 1.9864 \\
Non assicurata, \texttt{01:27:18}& 112.95 &  98.9335 & {---}  & 3.8105 & {---}  \\
\bottomrule
\end{tabular}
\end{table}

Nel primo caso il trader, assicurato con copertura \texttt{HIGH}, ha pagato un
premio pari allo \num{0,36}\,\% del valore in output della transazione per
fronteggiare una perdita del \num{3,51}\,\% del valore stesso: il rimborso
ricevuto (\num{1,9864}~MEVI) ha coperto pressoché integralmente l'importo
sottratto. Nel secondo caso il trader, privo di copertura, ha subìto una
perdita del \num{3,85}\,\% del valore della transazione --- avvenuta solo
65~secondi dopo la prima --- senza alcuna possibilità di rimborso. Il premio
che avrebbe pagato sottoscrivendo la copertura sarebbe stato dell'ordine dello
\num{0,36}\,\% del valore, trattandosi di una transazione eseguita nello
stesso intervallo temporale e di entità comparabile.

\section{L'assicurazione come strumento di pianificazione finanziaria}
L'utilità della copertura non si esaurisce nel rimborso diretto della perdita:
il vero vantaggio competitivo risiede nella \textbf{prevedibilità del costo}.
Senza assicurazione, un trader che opera su un DEX soggetto ad attacchi
sandwich è esposto a perdite inattese e non pianificabili: su ogni swap esiste
una probabilità che il bot sottragga una quota del valore della transazione,
con un impatto variabile e imprevedibile da un'operazione all'altra. Con la
copertura \texttt{HIGH}, invece, il costo di ogni swap è \emph{noto ex ante}:
il trader paga un premio fisso --- nell'ordine dello \num{0,38}\,\% del volume
--- indipendentemente da quante volte verrà attaccato, e può incorporarlo nei
propri calcoli esattamente come farebbe con le commissioni di gas o le fee
dell'AMM. La perdita da sandwich, se si verifica, rimane a carico del pool
assicurativo; l'incertezza individuale viene redistribuita sull'intera
comunità degli assicurati, secondo il principio attuariale di mutualizzazione
del rischio.

\section{Due trader a confronto: prevedibilità contro incertezza}

Per rendere concreta questa differenza, immaginiamo due trader --- chiamiamoli
A e B --- che eseguono ciascuno dieci swap di importo analogo, pari al valore
medio osservato nella simulazione: \num{74,44}~MEVI. Trader~A sottoscrive la
copertura \texttt{HIGH} per ogni transazione; Trader~B opera senza alcuna
protezione.

Prima ancora di eseguire il primo swap, Trader~A conosce esattamente quanto
pagherà in premi: lo \num{0,38}\,\% per ogni transazione, pari a circa
\num{0,28}~MEVI per swap, per un totale di \num{2,83}~MEVI sull'intera
sequenza di dieci operazioni. Questo numero non cambierà, indipendentemente
da quante volte verrà attaccato: se il bot lo colpisce una, due o anche tre
volte, il protocollo rimborserà integralmente la perdita e il costo netto
rimarrà comunque \num{2,83}~MEVI. Trader~A può inserire questa cifra nel
proprio piano finanziario come qualunque altra voce di spesa operativa.

Trader~B non paga alcun premio. Con un tasso di attacco dell'\num{8,55}\,\%,
ci si potrebbe aspettare in media circa un attacco su dieci swap --- ma questa
è una probabilità aggregata sull'intera popolazione di transazioni, non una
garanzia di distribuzione uniforme tra i singoli trader. Il bot seleziona le
vittime osservando la mempool: il volume della transazione, il gas price, il
momento esatto di invio. Nulla impedisce che Trader~B venga preso di mira zero
volte (nessuna perdita) oppure due o tre volte di fila, con perdite cumulate
del tutto non pianificate. La Tabella~\ref{tab:scenari} illustra come il costo
effettivo di Trader~B possa variare a seconda di come si distribuiscono
gli attacchi.

\begin{table}[h]
\centering
\caption{Costo effettivo su 10 swap a seconda degli attacchi subiti
         (perdita media per attacco: \num{3,52}~MEVI).}
\label{tab:scenari}
\small
\begin{tabular}{cccc}
\toprule
\textbf{Attacchi subiti} & \textbf{Trader~A (assicurato)}
  & \textbf{Trader~B (non assicurato)} & \textbf{Differenza} \\
\midrule
0 & \num{2,83}~MEVI & \num{0,00}~MEVI & B risparmia \num{2,83}~MEVI \\
1 & \num{2,83}~MEVI & \num{3,52}~MEVI & B paga \num{0,69}~MEVI in più \\
2 & \num{2,83}~MEVI & \num{7,04}~MEVI & B paga \num{4,21}~MEVI in più \\
3 & \num{2,83}~MEVI & \num{10,56}~MEVI & B paga \num{7,73}~MEVI in più \\
\bottomrule
\end{tabular}
\end{table}

La tabella evidenzia il problema strutturale di operare senza copertura. Se
Trader~B è fortunato e non viene mai attaccato, risparmia i \num{2,83}~MEVI
di premi. Ma se viene colpito due volte consecutive --- uno scenario tutt'altro
che raro quando la probabilità di attacco sfiora il 10\% --- la perdita sale a
\num{7,04}~MEVI, oltre due volte e mezzo il costo fisso dell'assicurazione,
senza preavviso e senza possibilità di rimborso.

Il punto cruciale è che la probabilità di attacco è una misura statistica
aggregata, non una garanzia di distribuzione uniforme tra i trader. Due
operatori con lo stesso profilo e lo stesso volume possono vivere esperienze
radicalmente diverse nello stesso arco temporale: uno esce illeso, l'altro
subisce due attacchi in pochi minuti. Non esiste
alcun meccanismo che garantisca equità nella distribuzione degli attacchi: il
bot si comporta in modo opportunistico, non casuale. L'assicurazione non
cambia la probabilità che un attacco avvenga, ma elimina l'incertezza su quanto
quel rischio costerà al singolo trader, convertendo l'imponderabile in un costo
fisso che può essere pianificato, budgetizzato e confrontato con qualunque altra
voce di spesa operativa.

\section{Confronto con Nexus Mutual}

Nexus Mutual è il principale protocollo di assicurazione decentralizzata su
Ethereum. Lanciato nel 2019, ha raccolto storicamente oltre trentasette milioni
di dollari in premi e ne ha distribuiti circa diciotto e mezzo in rimborsi,
raggiungendo un \emph{loss ratio} aggregato del \num{50,26}\,\%.\footnote{Dati
finanziari estratti dal dashboard ufficiale Nexus Mutual su Dune Analytics,
consultato maggio~2026: premi totali \$37.070.492, rimborsi totali
\$18.631.511.} Sebbene copra una tipologia di rischio diversa --- il fallimento
o l'exploit di smart contract, piuttosto che l'estrazione MEV --- il confronto
con MEV Insurance è significativo perché i due protocolli condividono la stessa
struttura economica di fondo e affrontano i loro rispettivi mercati con
soluzioni tecniche costruite sulle medesime premesse.

Entrambi si basano su un pool di capitali collettivo, alimentato dai premi
degli assicurati, che viene utilizzato per erogare i rimborsi quando un
sinistro viene giudicato legittimo. In entrambi i casi i partecipanti al
protocollo hanno un incentivo economico alla correttezza: Nexus Mutual remunera
gli assessori che votano in linea col consenso e penalizza quelli che si
discostano in modo anomalo, MEV Insurance distribuisce reward agli oracle che si
allineano alla mediana del FraudScore e applica uno slashing a chi devia. Anche
la gestione della solvibilità segue logiche parallele: quando le riserve
scendono sotto una soglia critica, entrambi i protocolli reagiscono alzando il
costo della copertura --- Nexus Mutual attraverso il suo parametro di capitale
minimo (MCR)\footnote{Il MCR (\textit{Minimum Capital Requirement}) è calibrato
a un livello di confidenza del 99{,}5\,\%. Cfr.\
\url{https://docs.nexusmutual.io/protocol/capital-pool/mcr/}.} e il meccanismo
di pricing dinamico RAMM,\footnote{Nexus Mutual ha adottato il meccanismo RAMM
(\textit{Ratcheting AMM}) nel novembre~2023 in sostituzione del precedente
sistema basato sulla curva di legame. Cfr.\
\url{https://docs.nexusmutual.io/protocol/pricing/}.} MEV Insurance attraverso
il margine adattivo che scala il premio in funzione del Solvency Ratio del pool.

La differenza più evidente riguarda il modo in cui viene verificata la
legittimità di un sinistro. Nexus Mutual affida questa decisione a un processo
di governance umana: un comitato di assessori esamina il caso, e se il consenso
non è sufficientemente ampio il voto viene esteso all'intera comunità dei
possessori di NXM, con un processo che richiede tipicamente dai tre ai sette
giorni.\footnote{La procedura prevede un minimo di 72~ore per il comitato di
primo livello e fino a 7~giorni in caso di voto allargato alla comunità. Cfr.\
\url{https://docs.nexusmutual.io/protocol/claims-assessment/}.} Questo approccio
è necessario quando il rischio coperto è di natura discrezionale: determinare
se una perdita è causata da un exploit e se rientra nei termini della polizza
richiede giudizio umano. MEV Insurance opera in modo radicalmente diverso,
perché un attacco sandwich lascia una firma on-chain matematicamente
inequivocabile, con frontrun, swap della vittima e backrun nello stesso blocco e
nello stesso ordine di esecuzione. Non serve interpretazione; serve verifica. La
rete di oracle esegue un processo commit-reveal in pochi minuti e produce un
verdetto algoritmico basato sul FraudScore, riducendo i tempi di liquidazione da
giorni a minuti e abbassando significativamente i costi operativi.

Il divario nel \emph{loss ratio} --- \num{98}\,\% per MEV Insurance contro il
\num{50}\,\% di Nexus Mutual --- riflette esattamente la diversa natura del
rischio coperto. Nexus Mutual mantiene un margine di sicurezza vicino al
\num{50}\,\% perché i rischi che affronta sono caratterizzati da alta
imprevedibilità e potenziale impatto catastrofico: un singolo exploit può
svuotare protocolli da centinaia di milioni di dollari, e le riserve devono
essere abbondanti per garantire la solvibilità anche negli scenari più estremi.
MEV Insurance copre invece un rischio statisticamente stabile e ad alta
frequenza, per il quale la formula attuariale del premio riesce a calibrare con
precisione l'esposizione del pool: il margine del \num{2}\,\% è sufficiente
perché la distribuzione degli attacchi è prevedibile in aggregato, anche se
imprevedibile per il singolo trader. In entrambi i casi, il \emph{loss ratio}
rimane inferiore al \num{100}\,\%: condizione necessaria affinché le riserve
crescano nel tempo e il protocollo rimanga solvibile nel lungo periodo.

Il punto più rilevante ai fini di questo lavoro non riguarda però le differenze
meccanistiche tra i due protocolli, bensì ciò che Nexus Mutual dimostra
empiricamente: un'assicurazione decentralizzata può funzionare. In oltre cinque
anni di operatività il protocollo ha distribuito milioni di dollari in rimborsi,
ha continuato a crescere in termini di capitale investito e di coperture
vendute,\footnote{Al 2025 Nexus Mutual gestisce un pool di capitali di circa
\$400~milioni e ha venduto oltre un miliardo di dollari di coperture nell'anno.
Fonte: Nexus Mutual, \textit{The Nexus Review --- Q4 2025 Wrap-Up},
\url{https://forum.nexusmutual.io/t/the-nexus-review-q4-2025-wrap-up/1824},
gennaio~2026.} e ha mantenuto un \emph{loss ratio} stabilmente inferiore al
\num{100}\,\%. Questa è la prova che gli utenti DeFi sono disposti a pagare
premi per proteggersi, che un meccanismo distribuito di verifica dei sinistri
può operare equamente su scala, e che un pool assicurativo decentralizzato può
rimanere solvibile nel lungo periodo. MEV Insurance si costruisce esattamente
su questi stessi presupposti e li applica a una classe di rischio più granulare.

Se Nexus Mutual ha dimostrato che un \emph{loss ratio} del \num{50}\,\% è
compatibile con un protocollo assicurativo credibile e in crescita, a maggior
ragione un \emph{loss ratio} del \num{98}\,\% --- più vicino alla neutralità
attuariale e quindi strutturalmente più vantaggioso per l'assicurato --- è
compatibile con la sostenibilità del pool, a condizione che la frequenza e la
prevedibilità degli eventi coperti siano sufficientemente stabili. La
simulazione dimostra che questa condizione è soddisfatta: il pool assicurativo è
cresciuto da \num{50} a \num{80,62}~MEVI nonostante il \emph{loss ratio} quasi
unitario, confermando che il margine del \num{2}\,\%, pur ridotto, è calibrato
correttamente sulla distribuzione statistica degli attacchi sandwich. Ciò che in
Nexus Mutual sarebbe un segnale di riserve insufficienti, nel contesto di MEV
Insurance è l'indicatore che i premi sono calibrati in modo quasi attuarialmente
equo: il protocollo non estrae valore in eccesso dagli assicurati, ma si limita
a coprire i propri costi operativi, esattamente come un'assicurazione matura e
ben governata dovrebbe fare.


\section{Considerazioni conclusive}
La simulazione, condotta con un tasso di attacco configurato al 10\,\% (tasso
effettivo: \num{8,55}\,\%) e una perdita media per attacco del \num{4,73}\,\%
del valore della singola transazione, conferma che la copertura \texttt{HIGH}
risulta economicamente conveniente per entrambe le parti: il trader è protetto
dalla perdita inattesa pagando un premio pari allo \num{0,38}\,\% del volume,
mentre il protocollo registra un \emph{loss ratio} del \num{98,00}\,\% --- un
margine del \num{2}\,\% sufficiente a coprire i costi operativi e ad accrescere
la riserva del pool da \num{50}~MEVI a \num{80,62}~MEVI nel corso della
simulazione.

L'esempio dei due trader (sezione~8) mostra come la convenienza
dell'assicurazione non risieda soltanto nel valore atteso dei rimborsi, ma
nella certezza del costo: il trader assicurato sa già quanto spenderà, quello
non assicurato non può prevederlo, e basta una sequenza di due attacchi
ravvicinati per trovarsi a sostenere perdite più che doppie rispetto al costo
fisso di una polizza. Il confronto con Nexus Mutual (sezione~9) evidenzia come
MEV Insurance si inserisca coerentemente nell'ecosistema delle assicurazioni
DeFi consolidate, condividendone i principi fondamentali e declinandoli in
forma parametrica e automatizzata, ideale per eventi con firma on-chain
deterministica come gli attacchi sandwich.

Il meccanismo di verifica degli oracoli, che in questa simulazione ha respinto
circa l'\num{8,26}\,\% dei \emph{claim}, si conferma essenziale per proteggere
il pool da rimborsi impropri. Il dato suggerisce che i client devono
implementare con precisione la costruzione del pattern sandwich affinché la
quasi totalità delle perdite legittime venga effettivamente indennizzata.

\end{document}
